\chapter{Mantenimiento y desarrollo}\label{cap:mantenimiento}

Este capítulo orienta a quien modifica el sistema: dónde está cada cosa, qué
reglas debe cumplir un cambio, cómo se verifica antes de integrarlo y qué
comprueba la integración continua. Las reglas completas están en
\archivo{CONTRIBUTING.md} y \archivo{AGENTS.md}; ambos archivos contienen las
mismas normas y deben mantenerse consistentes.

\section{Estructura del repositorio}\label{sec:mant-estructura}

\begin{table}[H]
  \centering
  \caption{Directorios y archivos principales del repositorio.}\label{tab:mant-estructura}
  \begin{tabularx}{\textwidth}{L{4.4cm} Y}
    \toprule
    \textbf{Ruta} & \textbf{Contenido} \\
    \midrule
    \archivo{crates/}
      & Espacio de trabajo de Rust. \archivo{crates/bin/} es la raíz de
        composición y produce \texttt{despacho-cli} y \texttt{qadra-ca}. \\
    \archivo{migrations/}
      & Migraciones SQL de PostgreSQL, numeradas \texttt{NNNN\_nombre.sql}. \\
    \archivo{web/}
      & Espacio de trabajo npm: \archivo{apps/qadra} (interfaz del despacho),
        \archivo{apps/ca} (portal de la CA) y \archivo{packages/ui}. \\
    \archivo{pki/}
      & Guiones y configuración de OpenSSL para la CA y la TSA internas. \\
    \archivo{ops/}
      & Automatización en Python del despliegue y de la operación de la CA. \\
    \archivo{scripts/}
      & Construcción de módulos WASM y de la versión, demostraciones,
        utilidades de CI y sus pruebas (\archivo{scripts/tests/}). \\
    \archivo{docs/}
      & Documentación técnica (índice en \archivo{docs/README.md}) y
        decisiones de arquitectura en \archivo{docs/adr/}. \\
    \archivo{latex/}, \archivo{manual-usuario/}, \archivo{manual-tecnico/}
      & Fuentes del reporte técnico y de los dos manuales. \\
    \archivo{.github/workflows/}
      & Flujos de integración continua y de despliegue. \\
    \archivo{Cargo.toml}, \archivo{deny.toml}
      & Manifiesto del espacio de trabajo y política de dependencias. \\
    \bottomrule
  \end{tabularx}
\end{table}

El grafo de crates impone la dirección de dependencias descrita en el
\cref{cap:arquitectura}; un cambio que agregue una dependencia en sentido
contrario no se integra.

\section{Reglas de contribución}\label{sec:mant-reglas}

\begin{itemize}
  \item \textbf{Solo ASCII} en código, comentarios, identificadores, nombres
    de archivo, configuración y mensajes de \emph{commit}. La prosa en
    Markdown y \LaTeX{} está exenta.
  \item \textbf{Archivos de menos de 400 líneas} de código en Rust, guiones y
    configuración con lógica. El módulo se divide antes de superar el límite.
  \item \textbf{Desarrollo guiado por pruebas}: primero la prueba que falla y
    después el código, ambos en el mismo cambio.
  \item \textbf{Comentarios autocontenidos}: no se citan planes, fases,
    tiques ni conversaciones. La justificación de una decisión se escribe en
    un registro de decisión de arquitectura (ADR) y se cita por su ruta.
  \item \textbf{La documentación viaja con el cambio}: el mismo
    \emph{pull request} actualiza \archivo{docs/http-api.md} y
    \archivo{docs/database-operations.md} cuando cambian rutas, permisos,
    migraciones u operación.
\end{itemize}

Los mensajes de \emph{commit} siguen \emph{Conventional Commits}
(\texttt{type(scope): summary}), con el resumen en imperativo, en minúsculas
y sin punto final. Los tipos admitidos son \texttt{feat}, \texttt{fix}, \texttt{docs},
\texttt{style}, \texttt{refactor}, \texttt{perf}, \texttt{test},
\texttt{build}, \texttt{ci}, \texttt{chore} y \texttt{revert}. Cada
\emph{commit} contiene un solo cambio lógico y se prepara indicando rutas
explícitas; no se usa \texttt{git add -A} ni \texttt{git add .}. No se trabaja
sobre \texttt{main}: las ramas se llaman \texttt{feat/<descripción>}
(implementación) o \texttt{progress/<descripción>} (documentación y
coordinación) y se integran mediante \emph{pull request} con
\emph{squash merge}.

Los registros de decisión (ADR) están en \archivo{docs/adr/} con el nombre
\texttt{NNNN-titulo-corto.md} y las secciones Contexto, Decisión, Estado y
Consecuencias. Cuando una decisión cambia, se conserva el registro anterior y
el nuevo identifica al que sustituye o enmienda.

\section{Verificación local antes de integrar}\label{sec:mant-verificacion}

Se ejecuta una sola campaña de verificación a la vez: no se solapan
compilaciones, pruebas, cobertura ni compilación de documentos. Si
\archivo{/tmp} es un sistema de archivos en memoria, el entorno se prepara
desde la raíz del repositorio:

\begin{verbatim}
mkdir -p output/tmp
chmod 700 output/tmp
export TMPDIR="$PWD/output/tmp"
export CARGO_BUILD_JOBS=1
export RUST_TEST_THREADS=1
\end{verbatim}

\subsection{Rust}\label{sec:mant-rust}

Las pruebas nativas de formatos documentales requieren Linux x86\_64 y la
biblioteca qpdf fijada, que se prepara antes de las pruebas. El cierre de una
entrega exige estas comprobaciones:

\begin{verbatim}
export TT_TEST_QPDF_LIBRARY="$(bash scripts/setup-document-formats.sh)"
cargo fmt --all
cargo build --workspace
cargo test --workspace
cargo clippy --workspace --all-targets -- -D warnings
cargo deny check
\end{verbatim}

\texttt{cargo deny check} aplica \archivo{deny.toml} (avisos de seguridad,
licencias, dependencias prohibidas y orígenes) y requiere tener instalado
\texttt{cargo-deny}. La versión mínima de Rust (1.88) se declara en
\archivo{Cargo.toml} y no se fija una cadena de herramientas local; la
integración continua ejecuta \texttt{clippy} con la versión estable vigente.

\importante{Las pruebas de los adaptadores de PostgreSQL y Redis terminan sin
ejecutarse cuando faltan sus variables de entorno
(\texttt{IDENTITY\_TEST\_DATABASE\_URL}, \texttt{IDENTITY\_TEST\_REDIS\_URL},
\texttt{CASE\_TEST\_DATABASE\_URL} y \texttt{DOCUMENT\_TEST\_DATABASE\_URL}).
Un resultado en verde de \texttt{cargo test} no prueba que se ejercitaron.
Para ejecutarlas use \texttt{bash scripts/test-backends.sh}, que crea bases de
datos aisladas y un Redis desechable y corre la suite completa.}

\subsection{Interfaz web y pruebas de navegador}\label{sec:mant-web}

Los dos módulos WASM se generan una vez desde la raíz y de nuevo cuando
cambian sus fuentes en Rust; sin ellos, \texttt{npm run dev} y
\texttt{npm run build} se detienen.

\begin{verbatim}
bash scripts/build-browser-crypto.sh
bash scripts/build-evidence-verify.sh
cd web
npm ci
npm run format:check
npm test
npm run test:ca
npm run build
npx playwright install chromium
npm run test:e2e -- --workers=1
\end{verbatim}

\texttt{npm test} y \texttt{npm run test:ca} ejecutan las pruebas unitarias
de Node de cada portal; \texttt{npm run test:e2e}, las de navegador con
servicios simulados (Playwright). Las pruebas de navegador contra servicios
reales y desechables (servidor Rust, PostgreSQL, Redis y TSA local) se lanzan
desde la raíz con \texttt{bash scripts/web-demo.sh}. Para cambios en la CLI se
ejecuta además \texttt{bash scripts/demo.sh} y, para cambios en el flujo HTTP
integrado, \texttt{bash scripts/api-demo.sh}.

\section{Integración continua}\label{sec:mant-ci}

Los flujos de \archivo{.github/workflows/} se ejecutan en \emph{runners}
propios del repositorio; su preparación está en
\archivo{docs/ci-runner-operations.md}. Cuando una comprobación falla, se
cancelan los trabajos restantes de la misma revisión.

\begin{table}[H]
  \centering
  \caption{Flujos de integración continua.}\label{tab:mant-flujos}
  \begin{tabularx}{\textwidth}{L{2.7cm} L{4.3cm} Y}
    \toprule
    \textbf{Flujo} & \textbf{Cuándo se ejecuta} & \textbf{Qué comprueba} \\
    \midrule
    \archivo{ci.yml} (CI)
      & \emph{Pull requests} que tocan algo fuera de \archivo{web/}, cada
        \emph{push} a \texttt{main}, una vez cada noche y a petición.
      & Formato, \texttt{clippy} sin avisos, pruebas de Rust y de los
        guiones, respaldo de despliegue, compilación con Rust 1.88,
        \texttt{cargo deny check}, tamaño del binario \texttt{despacho-cli}
        (máximo 25~MB) y cobertura. \\
    \archivo{web.yml} (Web)
      & \emph{Pull requests} y \emph{push} a \texttt{main} que tocan
        \archivo{web/}, \archivo{crates/}, \archivo{migrations/},
        \archivo{pki/} o los guiones que usa el flujo.
      & Módulos WASM, formato, pruebas unitarias, compilación, pruebas de
        navegador simuladas (dos particiones) y con servicios reales (tres
        particiones). \\
    \archivo{documents.yml} (Documents)
      & Cambios en \archivo{latex/}, publicación de una versión y a
        petición.
      & Compila el reporte con \texttt{make -B -C latex}, lo conserva como
        artefacto \texttt{report-pdf} y, en una versión publicada, lo adjunta
        con su \texttt{SHA256SUMS}. \\
    \archivo{deploy.yml} (Deploy version)
      & Etiquetas \texttt{v*} y a petición.
      & Reutiliza o repite las comprobaciones de CI y Web para el
        \emph{commit}, empaqueta con \archivo{scripts/build-release.sh} y,
        solo ante una etiqueta, despliega (\cref{cap:despliegue}). \\
    \bottomrule
  \end{tabularx}
\end{table}

\archivo{docs/adr/0088-ci-time-budget.md} fija un presupuesto de diez minutos
por campaña de \emph{pull request}. Por ello, un \emph{pull request} que solo
cambia \archivo{web/} no ejecuta CI, y el \emph{push} a \texttt{main} de un
árbol ya probado en su \emph{pull request} reutiliza ese resultado en lugar de
repetir las suites.

\subsection{Cobertura nocturna y umbral}\label{sec:mant-cobertura}

En el \emph{runner} dedicado, las pruebas de Rust de \emph{pull requests} y de
\texttt{main} se ejecutan sin instrumentación. La cobertura se mide cada noche
sobre \texttt{main} y cuando se lanza CI a mano con la entrada
\texttt{coverage} activada. \archivo{scripts/coverage-gate.sh} exige al menos
\textbf{90\,\%} de líneas cubiertas en cada uno de los crates \texttt{domain}, \texttt{application} e
\texttt{infrastructure}. El crate \texttt{bin} se informa, pero no tiene
umbral.

\importante{La cobertura no bloquea la integración ni el despliegue: una
caída por debajo del umbral aparece en la siguiente ejecución nocturna. Si un
\emph{pull request} elimina pruebas, lance CI con \texttt{coverage} activada
antes de integrarlo.}

\section{Cómo agregar una migración}\label{sec:mant-migracion}

Las migraciones no se leen del disco al ejecutar: el binario las incorpora al
compilar y las aplica en un orden fijo. Agregar una requiere cambios en SQL y
en Rust.

\begin{enumerate}
  \item Escriba primero la prueba que falla para el comportamiento nuevo.
  \item Cree \archivo{migrations/NNNN_nombre.sql} con el siguiente número
    libre. Los archivos de una misma entrega comparten número y se
    distinguen por el nombre, como \archivo{0052_contracts.sql} y
    \archivo{0052_contract_events.sql}.
  \item Registre el archivo con \texttt{include\_str!} en la constante
    \texttt{MIGRATIONS} del módulo de esquema correspondiente en
    \archivo{crates/infrastructure/src/}, por ejemplo
    \archivo{contract_schema/mod.rs}. Si el módulo es nuevo, agregue su
    recorrido en \archivo{crates/infrastructure/src/postgres.rs}, que fija
    el orden de aplicación.
  \item Actualice en ese módulo la validación del catálogo
    (\texttt{validate\_schema}) y los permisos del rol de ejecución
    (\texttt{grant\_runtime} y \texttt{validate\_runtime}, en
    \archivo{contract_schema/permissions.rs} para el ejemplo). El servidor
    valida el esquema al arrancar y rechaza uno incompleto o alterado.
  \item Ejecute \texttt{bash scripts/test-backends.sh} y aplique la
    migración en una base de datos de prueba con una
    \texttt{DATABASE\_URL} administrativa:
\begin{verbatim}
cargo build --workspace
target/debug/despacho-cli database migrate \
  --runtime-role tt_runtime
\end{verbatim}
  \item Documente el cambio y su procedimiento de aplicación en
    \archivo{docs/database-operations.md}.
\end{enumerate}

En el ejemplo, \texttt{tt\_runtime} es el nombre del rol de ejecución de
PostgreSQL; el comando no crea usuarios. La aplicación de una migración en un
entorno desplegado, con escritores detenidos y respaldo previo, se describe
en los \cref{cap:despliegue,cap:operacion}.

\section{Compilación de los documentos}\label{sec:mant-documentos}

Los tres documentos se compilan con LuaLaTeX y requieren TeX Live con las
fuentes Times New Roman y DejaVu Sans Mono. El reporte necesita además
\texttt{biber} y \texttt{makeglossaries}. Los PDF generados no se versionan.

\begin{verbatim}
make -C latex                 # reporte: latex/main.pdf
cd manual-usuario             # o manual-tecnico
latexmk -lualatex -outdir=build main.tex    # build/main.pdf
\end{verbatim}

Los detalles están en \archivo{latex/README.md},
\archivo{manual-usuario/README.md} y \archivo{manual-tecnico/README.md}. El flujo
Documents solo compila y publica el reporte: los dos manuales se compilan en
local. Después de compilar, revise en el PDF la paginación, las referencias
sin resolver y los glifos faltantes; las comprobaciones del reporte se
registran en \archivo{docs/academic-report-verification.md}.
